\chapter*{Заключение}
\addcontentsline{toc}{chapter}{Заключение}

В результате выполнения курсовой работы была разработана и исследована модель оценки возраста человека по изображению лица на основе сверточной нейронной сети. Для решения поставленной задачи была использована архитектура ResNet-50 с предобученными весами ImageNet, адаптированная под задачу регрессии.

В рамках работы были выполнены следующие подзадачи:
\begin{enumerate}
	\item проведен сбор и предобработка фотографий лиц людей;
	\item реализована модель на основе архитектуры ResNet-50 с заменой классификационной части на регрессионную голову;
	\item выполнено обучение модели с использованием обучающей и валидационной выборок с контролем качества по метрике MAE;
	\item выполнена финальная оценка качества модели на тестовой выборке и анализ полученных результатво.
\end{enumerate}

По результатам экспериментов было достигнуто значение средней абсолютной ошибки MAE, равное приблизительно 5 годам, что соответствует типичному уровню точности для задач оценки возраста по изображению лица. Анализ ошибок по возрастным группам показал, что наименьшая ошибка достигается для младших возрастов, тогда как для пожилых групп ошибка возрастает, что связано с дисбалансом данных и высокой визуальной вариативностью внешности.

В дальнейшем развитие работы может быть связано с использованием более сбалансированных датасетов, применением специализированных архитектур для анализа лиц, а также использованием дополнительных методов регуляризации и улучшения качества обучения.